\subsection{Process-based Analysis: How Action Composition Explains Harness Effects}
\label{sec:macro_process}
\vspace{-4pt}

To understand \emph{why} the same harness improves one model while degrading another, we decompose every model-harness trajectory into six non-overlapping action categories:
\begin{itemize}[leftmargin=*, itemsep=0pt, topsep=0.5pt]
    \item \emph{Understanding}: identifies the goal, givens, constraints, answer choices, and required output format.
    \item \emph{Planning}: decomposes the task or selects the next solving strategy.
    \item \emph{Reasoning}: performs model-internal solving, including derivation, calculation, or comparison.
    \item \emph{Tool Use}: calls an external tool or incorporates tool feedback into the solution.
    \item \emph{Verification}: checks, diagnoses, corrects, confirms, or rejects an intermediate result.
    \item \emph{Finalization}: states the final answer, boxed result, decisive conclusion, or terminal limitation.
\end{itemize}

Each trajectory step is classified into exactly one action via a deterministic regex-based classifier with the priority order Finalization \(>\) Tool Use \(>\) Verification \(>\) Reasoning \(>\) Planning \(>\) Understanding.
We aggregate action frequency distributions per benchmark, model, harness, and outcome.
Tables~\ref{tab:action-gpqa}, \ref{tab:action-hle}, and \ref{tab:action-aime} report the full decomposition; below we synthesize four cross-cutting findings that hold across benchmarks.

\findingbox[find:action-planning-collapse]{
\textbf{Planning Collapse under Tool-Use Harnesses.}
Harnessed agents consistently allocate a smaller share of their action budget to explicit planning than the same model run without a harness.
The effect is strongest for \texttt{Qwen3.5-27B} and \texttt{Kimi-K2.6}; \texttt{Gemma-4-31B-IT} already plans sparingly under direct inference, so the reduction is muted.
}

Across all benchmarks, \texttt{Qwen3.5-27B} Direct allocates 12.8--15.9\% of actions to Planning on success trajectories. Under \texttt{OpenClaw}, \texttt{OpenCode}, and \texttt{ZeroClaw}, this falls to 3.0--11.4\% — a reduction of 3.6--9.9 percentage points.
\texttt{Kimi-K2.6} shows a similar pattern: Planning drops from 7.8--11.9\% (Direct) to 3.6--7.2\% (tool-using agents) on GPQA and AIME, a loss of 2--6 pp.
\texttt{Gemma-4-31B-IT} Direct already allocates only 3.5--6.7\% to Planning; under agents the share changes by at most 3.5 pp in either direction.
%
The Planning collapse in \texttt{Qwen3.5} and \texttt{Kimi-K2.6} is not compensated by a corresponding increase in Tool Use on success: the model appears to offload cognitive structure onto tool interaction, replacing explicit decomposition with implicit, tool-guided reasoning.

\findingbox[find:action-tool-inversion]{
\textbf{Tool-Reasoning Inversion on Failure.}
When tool-using agents fail, Tool Use cannibalises Reasoning.
This inversion is severe for \texttt{Gemma-4-31B-IT} and \texttt{Kimi-K2.6} but largely absent for \texttt{Qwen3.5-27B}.
}

For \texttt{Gemma-4-31B-IT}, the most extreme case occurs on AIME under \texttt{OpenClaw}: Reasoning falls from 77.7\% (success) to 22.6\% (failure), while Tool Use surges from 9.7\% to 61.3\% — a 52 pp inversion.
\texttt{Kimi-K2.6} on AIME \texttt{OpenClaw} failure allocates \emph{all} its actions (100\%) to Tool Use, with Reasoning at 0\%; on GPQA \texttt{OpenClaw}, Tool Use rises from 25.3\% to 71.0\% while Reasoning collapses from 56.3\% to 23.9\%.
HLE shows a similar but milder pattern: \texttt{Gemma-4} \texttt{OpenClaw} Reasoning drops 13 pp with a matching Tool Use gain.
%
In contrast, \texttt{Qwen3.5-27B} on failure \emph{reduces} Tool Use and increases Reasoning — the opposite of the entanglement pattern. For instance, on GPQA \texttt{OpenClaw}, Qwen3.5's Tool Use falls from 6.5\% (success) to 0.2\% (failure). This suggests \texttt{Qwen3.5-27B} abandons tool strategies when they fail, while \texttt{Gemma-4} and \texttt{Kimi-K2.6} re-invoke tools, entering a tool-calling spiral.

\findingbox[find:action-verification-loop]{
\textbf{Verification Amplification in ZeroClaw.}
\texttt{ZeroClaw}, which reasons without external tools, amplifies Verification actions relative to Direct across all models — and this amplification is stronger on failure than on success.
}

On GPQA, Verification under \texttt{ZeroClaw} increases from 0.7\% (Gemma-4 Direct success) to 6.6\% (ZeroClaw success) and 26.3\% (ZeroClaw failure). For \texttt{Kimi-K2.6}, the corresponding values are 0.9\% \(\rightarrow\) 10.3\% \(\rightarrow\) 23.0\%. On HLE, the pattern persists: Gemma-4 Verification rises from 2.6\% (Direct) to 23.5\% (ZeroClaw success) and 17.8\% (ZeroClaw failure).
\texttt{Qwen3.5-27B} shows a weaker but consistent elevation: GPQA Verification moves from 8.8\% (Direct) to 7.1--11.7\% (ZeroClaw).
%
Importantly, Verification does not lead to Finalization. The chord diagrams confirm that Verification\(\leftrightarrow\)Reasoning forms the dominant bidirectional chord pair in ZeroClaw failure panels — the model checks, re-reasons, checks again, and fails to deliver an answer.

\findingbox[find:action-model-asymmetry]{
\textbf{Model-Harness Action Signatures.}
The three models exhibit distinct action-composition signatures that explain why the same harness produces opposite outcome effects.
}

We identify three qualitatively different regimes:
\begin{itemize}[leftmargin=*, itemsep=0pt, topsep=0.5pt]
    \item \textbf{\texttt{Qwen3.5-27B}: Solve-dominant, tool-averse.}
    Reasoning occupies 69--83\% of all actions across settings. Under Direct, the model already allocates 9--14\% to Verification. Harnesses reduce Planning without adding Tool Use; failure trajectories show \emph{increased} Reasoning and \emph{decreased} Tool Use. The model retreats to internal reasoning when tools fail, but loses the planning structure needed to organize that reasoning — producing long, unguided reasoning chains that rarely reach Finalization.

    \item \textbf{\texttt{Gemma-4-31B-IT}: Tool-entangled, verification-prone.}
    Under Direct, Gemma-4 allocates 8--12\% to Finalization (vs.\ <1\% for Qwen3.5) and only 1--3\% to Verification. Harnesses increase both Tool Use (to 6--25\%) and, in ZeroClaw, Verification (to 7--26\%). On failure, Tool Use explodes and Reasoning collapses — the model is `trapped' by its own tool calls. This pattern is consistent with the outcome data (Tab.~\ref{tab:macro-main-results}), where Gemma-4 benefits from harnesses on computation-heavy tasks (AIME, GPQA) but the tool entanglement penalty appears mainly in failure trajectories.

    \item \textbf{\texttt{Kimi-K2.6}: Bimodal, high-variance.}
    Kimi-K2.6 shows the widest spread. Under Direct on HLE, it allocates 53\% to Finalization — far more than Qwen3.5 (<1\%) or Gemma-4 (12\%). Under agents, Finalization collapses to 5--12\%. On AIME \texttt{OpenClaw} failure, 100\% of actions are Tool Use. The model oscillates between efficient tool use (producing strong GPQA and AIME gains under ZeroClaw/OpenCode) and catastrophic tool spirals (producing the 46 pp GPQA drop under OpenClaw). Its high Direct Finalization rate suggests strong internal answer-generation capability that harnesses disrupt rather than augment.
\end{itemize}
